% Part: first-order-logic
% Chapter: syntax-and-semantics
% Section: subformulas

\documentclass[../../../include/open-logic-section]{subfiles}

\begin{document}

\olfileid[tr]{fol}{syn}{sbf}

\olsection{\printtoken{P}{subformula}}

\begin{explain}
Belirli bir formülü ``oluşturan'' !!{formula}s hakkında konuşmak çoğu
kez yararlıdır. Verilen !!{formula} için bunları
\emph{!!{subformula}s} diye adlandırırız. Her !!{formula}, kendi kendisinin bir
!!a{subformula} sayılır; $!A$ söz konusu olduğunda, $!A$'nın kendisi dışındaki bir alt formüle
\emph{öz !!{subformula}} denir.
\end{explain}

\begin{defn}[Dolaysız !!^{subformula}]
$!A$ bir !!a{formula} ise \emph{dolaysız !!{subformula}s}, $!A$ için
aşağıdaki gibi tümevarımsal olarak tanımlanır:
\begin{enumerate}
\item Atomik !!{formula}s için dolaysız !!{subformula}s yoktur.

\tagitem{prvNot}{\indcase{!A}{\lnot !B}{$\indfrm$'nin tek dolaysız
    !!{subformula}{}ü $!B$'dir.}}{}

\item \indcase{!A}{(!B \ast !C)}{$\indfrm$'nin dolaysız
  !!{subformula}s{}i $!B$ ve $!C$'dir ($\ast$, iki yerli bağlaçlardan
  herhangi biridir).}

\tagitem{prvAll}{\indcase{!A}{\lforall[x][!B]}{$\indfrm$'nin tek dolaysız
    !!{subformula}{}ü $!B$'dir.}}{}

\tagitem{prvEx}{\indcase{!A}{\lexists[x][!B]}{$\indfrm$'nin tek dolaysız
    !!{subformula}{}ü $!B$'dir.}}{}
\end{enumerate}
\end{defn}

\begin{defn}[Öz !!^{subformula}]
$!A$ bir !!a{formula} ise \emph{öz !!{subformula}s}, $!A$ için
aşağıdaki gibi özyinelemeli olarak tanımlanır:
\begin{enumerate}
\item Atomik !!{formula}s için öz !!{subformula}s yoktur.

\tagitem{prvNot}{\indcase{!A}{\lnot !B}{$\indfrm$'nin öz
    !!{subformula}s{}i, $!B$ ile $!B$'nin bütün öz !!{subformula}s{}idir.}}{}

\item \indcase{!A}{(!B \ast !C)}{$\indfrm$'nin öz !!{subformula}s{}i;
  $!B$, $!C$, $!B$'nin bütün öz !!{subformula}s{}i ve $!C$'nin bütün öz
  alt formülleridir.}

\tagitem{prvAll}{\indcase{!A}{\lforall[x][!B]}{$\indfrm$'nin öz
    !!{subformula}s{}i, $!B$ ile $!B$'nin bütün öz !!{subformula}s{}idir.}}{}

\tagitem{prvEx}{\indcase{!A}{\lexists[x][!B]}{$\indfrm$'nin öz
    !!{subformula}s{}i, $!B$ ile $!B$'nin bütün öz !!{subformula}s{}idir.}}{}
\end{enumerate}
\end{defn}

\begin{defn}[!!^{subformula}]
$!A$'nın !!{subformula}s{}i, $!A$'nın kendisi ile bütün öz
!!{subformula}s{}dir.
\end{defn}

\begin{explain}
Dolaysız !!{subformula}s ile öz !!{subformula}s arasındaki tanım
farkına dikkat edin. İlkinde, bir formül~$!A$ için, $!A$'nın her olası
biçimindeki dolaysız !!{subformula}s kümesini doğrudan tanımladık. Bu,
durumlara göre açık bir tanımdır ve durumlar
!!{formula}s kümesinin tümevarımsal tanımını yansıtır. İkincisinde yine
bütün !!{formula}s kümesinin tanımlanma biçimini yansıttık; ancak her
durumda öz !!{subformula}s kümesini de kullandık: daha küçük
!!{formula}s $!B$ ile $!C$ için bu kümeyi, söz konusu !!{formula}s ile
birlikte tanıma kattık. Bu, tanımı \emph{özyinelemeli}
kılar. Genel olarak, tümevarımsal tanımlanmış bir küme (bizim
durumumuzda !!{formula}s) üzerindeki bir fonksiyonun tanımı, fonksiyon
tanımındaki durumlar fonksiyonun kendisini kullanıyorsa özyinelemelidir.
Ancak tanımın iyi tanımlanmış olması için yalnızca tanımladığımız
argümandan ``önce'' gelen argümanlardaki fonksiyon değerlerini
kullandığımızdan emin olmalıyız---bizim durumumuzda $(!B \ast !C)$ için
``öz !!{subformula}'' tanımlarken yalnızca öz !!{subformula}s kümesini, yani
``daha önceki'' !!{formula}s $!B$ ile $!C$ için olanları kullanırız.
\end{explain}

\begin{prop}
\ollabel{prop:subfrm-trans}
$!B$, $!A$'nın bir alt formülü ve $!C$, $!B$'nin bir alt formülü olsun.
O zaman $!C$, $!A$'nın bir alt formülüdür. Başka bir deyişle alt formül
bağıntısı geçişlidir.
\end{prop}

\begin{prob}
\olref[fol][syn][sbf]{prop:subfrm-trans}'i ispatlayın.
\end{prob}

\begin{prop}
\ollabel{prop:count-subfrms}
$!A$, $n$ bağlaç ve niceleyici içeren bir formül olsun. O zaman $!A$'nın
en çok $2n+1$ alt formülü vardır.
\end{prop}

\begin{prob}
\olref[fol][syn][sbf]{prop:count-subfrms}'i ispatlayın.
\end{prob}

\end{document}
